\chapter{Fields for Free Particles and Discrete Symmetries}

\section{Overview: from relativistic wave equations to QFT}
We need a formalism that:
\begin{itemize}
  \item combines Quantum Mechanics (QM) and Special Relativity (SR),
  \item describes processes with an arbitrary number of particles.
\end{itemize}
The outcome is \emph{Quantum Field Theory} (QFT). A systematic way to build a QFT for free particles is:
\begin{enumerate}
  \item Find a suitable (relativistic) wave equation.
  \item Find a Lagrangian whose Euler--Lagrange equations reproduce it.
  \item Apply canonical quantization rules.
\end{enumerate}
In QFT, wave functions become \emph{operators}, and the arbitrary coefficients in general solutions
become \emph{creation and annihilation operators}. We will go through:
\begin{itemize}
  \item Schr\"odinger field (non-relativistic),
  \item scalar field (spin 0, Klein--Gordon),
  \item Dirac field (spin $1/2$),
  \item vector field (spin 1, Proca/Maxwell),
\end{itemize}
and we will discuss the implementation of discrete symmetries P, C and T in each case.

\section{The Schr\"odinger field (non-relativistic)}
\subsection{Equation of motion and Lagrangian}
The Schr\"odinger equation is
\begin{equation}
\left(i\frac{\partial}{\partial t}+\frac{\nabla^2}{2m}\right)\psi(t,\vec x)=0.
\end{equation}
If the particle has spin $s$, each of the $2s+1$ components satisfies the same equation.
This equation follows from the action $S=\int dt\,L$, with
\begin{equation}
L=\int d^3x\,\mathcal L,\qquad
\mathcal L(t,\vec x)=\psi^\dagger(t,\vec x)\left(i\partial_0+\frac{\nabla^2}{2m}\right)\psi(t,\vec x).
\end{equation}

\subsection{Canonical quantization and Fock space}
Define the canonical momenta (for components $a=1,\dots,2s+1$)
\begin{equation}
\Pi_a(x)\equiv \frac{\partial\mathcal L}{\partial(\partial_0\psi_a(x))}= i\,\psi_a^\ast(x),
\qquad
\Pi_a^\ast(x)\equiv \frac{\partial\mathcal L}{\partial(\partial_0\psi_a^\ast(x))}=0.
\end{equation}
Canonical quantization at $t=0$ promotes fields to operators and imposes
\begin{equation}
[\hat\psi_a(\vec x),\hat\Pi_b(\vec y)]= i\,\delta^{(3)}(\vec x-\vec y)\,\delta_{ab},
\qquad
[\hat\psi_a(\vec x),\hat\psi_b(\vec y)]=[\hat\Pi_a(\vec x),\hat\Pi_b(\vec y)]=0,
\end{equation}
with the rule that for fermions commutators are replaced by anticommutators.

The general solution becomes an operator expansion
\begin{equation}
\hat\psi(\vec x,t)=\int\frac{d^3p}{(2\pi)^3}\sum_{\lambda=-s}^{s}
\hat a_\lambda(\vec p)\,\chi_\lambda\,e^{-iEt+i\vec p\cdot\vec x},\qquad
E=\frac{\vec p^{\,2}}{2m},
\end{equation}
with $[\hat a_\lambda(\vec p),\hat a^\dagger_{\lambda'}(\vec p\,')]=(2\pi)^3\delta^{(3)}(\vec p-\vec p\,')\delta_{\lambda\lambda'}$
(and commutators among $\hat a$ or among $\hat a^\dagger$ vanishing, for bosons).

Assuming a vacuum state $|0\rangle$ such that $\hat a_\lambda(\vec p)|0\rangle=0$, one defines
one-particle states $|\vec p\,\lambda\rangle=\hat a^\dagger_\lambda(\vec p)|0\rangle$.
Multi-particle states follow similarly, spanning the \emph{Fock space}.

\subsection{Discrete symmetries: parity and time reversal}
Parity acts as $\vec x\to -\vec x$ and $\vec p\to -\vec p$, with $\vec L\to \vec L$ and $\vec S\to \vec S$.
In QFT it is implemented by a unitary operator $P$ such that $P^2=1$ and
\begin{equation}
P\,\psi(t,\vec x)\,P^{-1}=\pm \psi(t,-\vec x),
\qquad
P\,a_\lambda(\vec p)\,P^{-1}=\pm a_\lambda(-\vec p).
\end{equation}
Time reversal acts as $\vec x\to \vec x$, $\vec p\to -\vec p$, $\vec L\to -\vec L$, $\vec S\to -\vec S$.
It is implemented by an \emph{antiunitary} operator $T$.

\subsection{Coupling to electromagnetism: gauge invariance and minimal coupling}
The Maxwell equations in vacuum follow from
\begin{equation}
\mathcal L_{\rm EM}=-\frac14 F_{\mu\nu}F^{\mu\nu},\qquad
F_{\mu\nu}=\partial_\mu A_\nu-\partial_\nu A_\mu,
\end{equation}
which is invariant under the gauge transformation $A_\mu\to A_\mu-\partial_\mu\theta(x)$.
For a charged Schr\"odinger field, gauge symmetry acts as $\psi(x)\to e^{iq\theta(x)}\psi(x)$.
Gauge invariance is ensured by \emph{minimal coupling}
\begin{equation}
\partial_\mu\to D_\mu\equiv \partial_\mu+i q A_\mu,
\end{equation}
leading schematically to $\mathcal L=\mathcal L_{\rm EM}+\psi^\dagger\left(iD_0+\frac{\vec D^{\,2}}{2m}+\cdots\right)\psi$.
If $s\neq 0$, non-minimal couplings (e.g.\ magnetic moment terms) may also exist.

\section{Scalar fields: Klein--Gordon (spin 0)}
\subsection{Wave equation and Lagrangian}
The simplest relativistic wave equation is the Klein--Gordon equation
\begin{equation}
(\partial_\mu\partial^\mu + m^2)\,\phi(x)=0.
\end{equation}
For a complex scalar field, it follows from
\begin{equation}
\mathcal L=(\partial_\mu\phi^\ast)(\partial^\mu\phi)-m^2\phi^\ast\phi.
\end{equation}
For a real scalar field, one uses
\begin{equation}
\mathcal L=\frac12(\partial_\mu\phi)(\partial^\mu\phi)-\frac12 m^2\phi^2.
\end{equation}

\subsection{Mode expansion and particle/antiparticle operators}
Upon quantization, the general solution can be written as
\begin{equation}
\hat\phi(x)=\int\frac{d^3p}{(2\pi)^3}\frac{1}{\sqrt{2E_{\vec p}}}
\left(e^{-ip\cdot x}\,\hat a(\vec p)+e^{ip\cdot x}\,\hat b^\dagger(\vec p)\right),
\qquad
E_{\vec p}=\sqrt{\vec p^{\,2}+m^2}.
\end{equation}
Here $\hat a$ annihilates a spin-0 particle and $\hat b$ annihilates its antiparticle. For a real scalar field,
particle and antiparticle coincide: $\hat a=\hat b$.

The vacuum $|0\rangle$ satisfies $\hat a(\vec p)|0\rangle=\hat b(\vec p)|0\rangle=0$.
With relativistic normalization, one often chooses factors so that
$\langle \vec p\,|\,\vec p\,'\rangle = 2E_{\vec p}(2\pi)^3\delta^{(3)}(\vec p-\vec p\,')$.

\subsection{Discrete symmetries P, C and T}
Parity can be implemented as
\begin{equation}
P\,\phi(t,\vec x)\,P^{-1}=\pm \phi(t,-\vec x),
\qquad
P\,a(\vec p)\,P^{-1}=\pm a(-\vec p),\quad
P\,b(\vec p)\,P^{-1}=\pm b(-\vec p),
\end{equation}
so particle and antiparticle have the \emph{same} intrinsic parity for scalars.

Charge conjugation exchanges particle and antiparticle operators:
\begin{equation}
C\,a(\vec p)\,C^{-1}=b(\vec p),\qquad C\,b(\vec p)\,C^{-1}=a(\vec p),
\end{equation}
and acts on the field as $C\,\phi(x)\,C^{-1}=\phi^\ast(x)$.
For a real scalar field, $C$ can act as $C\phi C^{-1}=\pm\phi$ (defining $C$-parity).

Time reversal can be implemented as
\begin{equation}
T\,\phi(t,\vec x)\,T^{-1}=\eta_T\,\phi(-t,\vec x),\qquad |\eta_T|=1.
\end{equation}

\subsection{Coupling to electromagnetism}
Assign the gauge transformation $\phi(x)\to e^{iq\theta(x)}\phi(x)$ and apply minimal coupling:
\begin{equation}
\partial_\mu\to D_\mu=\partial_\mu+i q A_\mu,\qquad
\mathcal L = -\frac14 F_{\mu\nu}F^{\mu\nu} + (D_\mu\phi)^\ast(D^\mu\phi)-m^2\phi^\ast\phi.
\end{equation}
Minimal coupling implies that particle and antiparticle carry opposite charges.
For a real scalar field, only non-minimal couplings are allowed.

\section{Dirac fields: spin $1/2$}
\subsection{Dirac equation and Lagrangian}
The relativistic wave equation for spin-$1/2$ particles is the Dirac equation
\begin{equation}
(i\gamma^\mu\partial_\mu - m)\,\psi(x)=0,
\end{equation}
where $\psi$ is a complex 4-component spinor and the gamma matrices satisfy
\begin{equation}
\{\gamma^\mu,\gamma^\nu\}=2g^{\mu\nu}.
\end{equation}
The corresponding Lagrangian is
\begin{equation}
\mathcal L=\bar\psi\,(i\slashed{\partial}-m)\psi,
\qquad
\bar\psi\equiv \psi^\dagger\gamma^0,\qquad
\slashed{\partial}\equiv \gamma^\mu\partial_\mu.
\end{equation}

\subsection{Mode expansion and quantization}
The quantized field admits the standard expansion
\begin{equation}
\hat\psi(x)=\int\frac{d^3p}{(2\pi)^3}\frac{1}{\sqrt{2E_{\vec p}}}
\sum_{\lambda=\pm}\left(
e^{-ip\cdot x}\,u_\lambda(\vec p)\,\hat a_\lambda(\vec p)
+
e^{ip\cdot x}\,v_\lambda(\vec p)\,\hat b^\dagger_\lambda(\vec p)
\right),
\qquad
E_{\vec p}=\sqrt{\vec p^{\,2}+m^2},
\end{equation}
where $u_\lambda$ and $v_\lambda$ satisfy
\begin{equation}
(\slashed{p}-m)u_\lambda(\vec p)=0,\qquad (\slashed{p}+m)v_\lambda(\vec p)=0,\qquad \slashed{p}\equiv \gamma^\mu p_\mu.
\end{equation}
Consistency requires \emph{fermionic} quantization: the creation/annihilation operators satisfy anticommutation relations.

\subsection{Gamma-matrix technology (useful identities)}
A complete basis of $4\times4$ matrices is $\{I,\gamma^5,\gamma^\mu,\gamma^5\gamma^\mu,\sigma^{\mu\nu}\}$ with
$\sigma^{\mu\nu}\equiv \frac{i}{2}[\gamma^\mu,\gamma^\nu]$.
Important trace properties include:
\begin{itemize}
  \item $\mathrm{tr}(\text{odd number of }\gamma)=0$,
  \item $\mathrm{tr}(\gamma^\mu\gamma^\nu)=4g^{\mu\nu}$,
  \item $\mathrm{tr}(\gamma^\mu\gamma^\nu\gamma^\rho\gamma^\sigma)
  =4(g^{\mu\nu}g^{\rho\sigma}-g^{\mu\rho}g^{\nu\sigma}+g^{\mu\sigma}g^{\nu\rho})$,
  \item $\mathrm{tr}(\gamma^5\gamma^\mu\gamma^\nu\gamma^\rho\gamma^\sigma)=4i\epsilon^{\mu\nu\rho\sigma}$.
\end{itemize}

\subsection{Discrete symmetries P, C and T}
A standard implementation is:
\begin{align}
P\,\psi(t,\vec x)\,P^{-1} &= \gamma^0\,\psi(t,-\vec x), \\
C\,\psi(x)\,C^{-1} &= i\gamma^2\,\psi^\ast(x)\equiv \psi^c(x), \\
T\,\psi(t,\vec x)\,T^{-1} &= \eta_T\,\gamma^1\gamma^3\,\psi(-t,\vec x),
\end{align}
with $\eta_T$ a phase. One consequence is that particle and antiparticle have \emph{opposite} parity for Dirac fields
(visible already at the level of how $a_\lambda$ and $b_\lambda$ transform).

\subsection{Coupling to electromagnetism}
Gauge invariance under $\psi\to e^{iq\theta(x)}\psi$ is ensured by minimal coupling:
\begin{equation}
\partial_\mu\to D_\mu=\partial_\mu+i q A_\mu,\qquad
\mathcal L = -\frac14 F_{\mu\nu}F^{\mu\nu}+\bar\psi(i\slashed{D}-m)\psi.
\end{equation}
Particle and antiparticle have opposite electric charge.

\section{Chirality, helicity and Weyl fermions}
Define the chiral projectors
\begin{equation}
P_R=\frac{1+\gamma^5}{2},\qquad P_L=\frac{1-\gamma^5}{2},\qquad
\psi_R\equiv P_R\psi,\qquad \psi_L\equiv P_L\psi,
\end{equation}
with $P_R+P_L=1$, $P_R^2=P_R$, $P_L^2=P_L$, and $P_RP_L=0$.

In the massless limit (equivalently, at very high energies $E\gg m$), right- and left-handed components decouple:
\begin{equation}
\bar\psi(i\slashed{\partial}-m)\psi
= \bar\psi_R i\slashed{\partial}\psi_R+\bar\psi_L i\slashed{\partial}\psi_L
- m\bar\psi_R\psi_L - m\bar\psi_L\psi_R
\;\;\xrightarrow[E\gg m]{}\;\;
\bar\psi_R i\slashed{\partial}\psi_R+\bar\psi_L i\slashed{\partial}\psi_L.
\end{equation}
In this regime, chirality aligns with helicity: right-handed fields describe particles with positive helicity
(and antiparticles with negative helicity), while left-handed fields describe particles with negative helicity
(and antiparticles with positive helicity). The two-component Weyl equations arise as the building blocks of the electroweak theory.

Discrete symmetries act non-trivially on chirality; in particular, parity exchanges left and right components.

\section{Majorana masses and Majorana fermions}
Besides the usual Dirac mass term, one can add a \emph{Majorana mass} term of the schematic form
\begin{equation}
\delta\mathcal L = -\delta m\,(\bar\psi\,\psi^c+\bar\psi^c\,\psi),
\end{equation}
which violates the global U(1) symmetry $\psi\to e^{i\theta}\psi$ and therefore does not conserve
``particle number minus antiparticle number''.

One can define a Majorana fermion by the condition
\begin{equation}
\psi=\psi^c=i\gamma^2\psi^\ast,
\end{equation}
implying that particle and antiparticle operators coincide (roughly $a_\lambda=b_\lambda$), so $\psi$ and $\psi^\ast$
are no longer independent.

\section{Vector fields: Proca and Maxwell (spin 1)}
\subsection{Equation of motion and Lagrangians}
A massive spin-1 field $B_\mu$ satisfies the Proca equation
\begin{equation}
\partial_\mu B^{\mu\nu}+m^2 B^\nu=0,\qquad B^{\mu\nu}=\partial^\mu B^\nu-\partial^\nu B^\mu.
\end{equation}
For $m=0$ this reduces to Maxwell's equations (in vacuum) for a gauge field.

For a \emph{real} vector field the Lagrangian is
\begin{equation}
\mathcal L = -\frac14 B_{\mu\nu}B^{\mu\nu}+\frac12 m^2 B_\mu B^\mu.
\end{equation}
For a \emph{complex} vector field,
\begin{equation}
\mathcal L = -\frac12 B_{\mu\nu}^\ast B^{\mu\nu}+ m^2 B_\mu^\ast B^\mu.
\end{equation}

\subsection{Mode expansion and polarizations}
For $m\neq 0$, the quantized field can be expanded as
\begin{equation}
B_\mu(x)=\int\frac{d^3p}{(2\pi)^3}\frac{1}{\sqrt{2E_{\vec p}}}
\sum_{h=0,\pm}\left(
e^{-ip\cdot x}\,\epsilon_\mu^{(h)}(\vec p)\,a_h(\vec p)
+
e^{ip\cdot x}\,\epsilon_\mu^{(h)\ast}(\vec p)\,a_h^\dagger(\vec p)
\right),
\end{equation}
with transversality $p^\mu \epsilon_\mu^{(h)}(\vec p)=0$ and completeness relations that project onto the physical
spin-1 subspace.

\subsection{Discrete symmetries and coupling to electromagnetism}
Parity, charge conjugation and time reversal can be implemented consistently; for complex fields,
charge conjugation exchanges particle and antiparticle operators $a_h\leftrightarrow b_h$.

If $B_\mu$ is complex and charged, gauge invariance under $B_\mu\to e^{iq\theta(x)}B_\mu$ is ensured by minimal coupling:
\begin{equation}
\partial_\mu \to D_\mu=\partial_\mu+i q A_\mu,\qquad
\tilde B_{\mu\nu}\equiv D_\mu B_\nu - D_\nu B_\mu,
\end{equation}
leading to a gauge-invariant Lagrangian of the form
\begin{equation}
\mathcal L = -\frac14 F_{\mu\nu}F^{\mu\nu} -\frac12 \tilde B_{\mu\nu}^\ast \tilde B^{\mu\nu} + m^2 B_\mu^\ast B^\mu.
\end{equation}
For a real vector field, only non-minimal couplings are allowed (unless the field itself is the gauge field).

\section{Summary}
The free-field construction illustrates the general QFT pattern:
\begin{itemize}
  \item Start from a wave equation, derive it from a Lagrangian,
  \item quantize canonically to obtain creation/annihilation operators and Fock space,
  \item implement discrete symmetries via unitary/antiunitary operators,
  \item include interactions by adding local terms that respect the symmetries (e.g.\ gauge invariance via minimal coupling).
\end{itemize}
This framework will be the foundation for interacting theories and for computing physical observables in the SM.